% Vince proofread

%========================================
% Section [12] Ovāda-pāṭimokkha-gāthā
%========================================
\section*{m12 | Ovāda-pāṭimokkha-gāthā}

\begin{chantsubsection}
  \chantnote
    {\{Handa mayaṃ buddhassa ovādanusāsaniyo bhaṇāmase\}}
  \chantnotetrans
    {\{Now let us recite the Buddha's words of advice and instruction.\}}
\end{chantsubsection}
\chantgap
\begin{chantsubsection}
  % 第一段 Khantī...
  \chantnote
    {Khantī paramaṃ tapo titikkhā,}
  \chantindent
    {Patience is the supreme purifying practice,}
  \par\chantindent
    {(khuam-ot-thon khue khuam-ot-klan pen ta-ba yang-ying)}
  \chantgap

  \chantnote
    {Nibbānaṃ paramaṃ vadanti Buddhā,}
  \chantindent
    {Nibbāna is supreme, say the Buddhas.}
    \par
  \chantindent
    {(phra-phut-tha-jao tang-lai trat-wa phra-nip-phan pen-yiam)}
  \chantgap

  \chantnote
    {Na hi pabbajito parūpaghātī, samaṇo hoti paraṃ viheṭhayanto.}
  \chantindent
    {A mendicant does not harm others, a recluse oppresses no one.}
  \par\chantindent
    {(phu-lang-phlan phoo-uen mai-chue wa-pen ban-pha-chit, phoo-biat-bian phoo-uen mai-chue-wa pen-samana-loei)}
  \chantgap

  \chantnote
    {Etaṃ Buddhānaṃ sāsanaṃ.}
  \chantindent
    {This is the teaching of all Buddhas.}
    \par
  \chantindent
    {(nī-pen kham-sorn khong phra-phut-tha-jao tang-lai)}
  \chantgap

  % 第二段 Sabbapāpassa...
  \chantnote
    {Sabbapāpassa akaraṇaṃ,}
  \chantindent
    {Abstaining from all evil,}
  \par\chantindent
    {(kan-mai-tham-bap-tang-puang nueng)}
  \chantgap

  \chantnote
    {kusalassūpasampadā,}
  \chantindent
    {cultivating the wholesome,}
    \par
  \chantindent
    {(kan-bam-phen ku-son hai-thueng-phrom nueng)}
  \chantgap

  \chantnote
    {sacittapariyodapanaṃ,}
  \chantindent
    {cleansing of one's own mind,}
    \par
  \chantindent
    {(kan-klan-chit khong-ton hai-phong-phaew nueng)}
  \chantgap

  \chantnote
    {etaṃ Buddhānaṃ sāsanaṃ.}
  \chantindent
    {This is the teaching of all Buddhas.}
  \par\chantindent
    {(nī pen-kham-sorn khong phra-phut-tha-jao tang-lai)}
  \chantgap

  \clearpage
  
  % 第三段 Anūpavādo...
  \chantnote
    {anūpavādo,}
  \chantindent
    {not insulting,}
  \par\chantindent
    {(kan-mai-khao-pai wa-rai-kan nueng)}
  \chantgap

  \chantnote
    {anūpaghāto,}
  \chantindent
    {not harming,}
  \par\chantindent
    {(kan-kai-khao-pai lang-phlan-kan nueng)}
  \chantgap

  \chantnote
    {pāṭimokkhe ca saṃvaro,}
  \chantindent
    {carefulness in the Fundamental Precepts,}
  \par\chantindent
    {(khuam-sam-ruam nai phra-pa-ti-mok nueng)}
  \chantgap

  \chantnote
    {mattaññutā ca bhattasmiṃ,}
  \chantindent
    {moderation in food,}
  \par\chantindent
    {(khuam-pen phoo-ru-pra-man nai pho-cha-na-han nueng)}
  \chantgap

  \chantnote
    {pantañca sayanāsanaṃ,}
  \chantindent
    {dwelling in solitude,}
  \par\chantindent
    {(kan-non kan-nang an-sa-ngad nueng)}
  \chantgap

  \chantnote
    {adhicitte ca āyogo,}
  \chantindent
    {engaging in higher mental development,}
  \par\chantindent
    {(kan-pra-kop khuam-phian nai a-thi-jit nueng)}
  \chantgap

  \chantnote
    {etaṃ Buddhānaṃ sāsanaṃ.}
  \chantindent
    {This is the teaching of all Buddhas.}
  \par\chantindent
    {(nī pen-kham-sorn khong phra-phut-tha-jao tang-lai)}
\end{chantsubsection}

    
% \end{chantsubsection}

% 原文（需要校对字符）
% ```
% Ovada-patimokkha-gatha
% {Handa mayam buddhassa ovadanusasanlyo bhanama se.}

% KhantI paramam tapo titikkha
% Nibbanam paramam vadanti buddha
% Na hi pabbajito parupaghati
% Samano hoti param vihethayanto
% Etam buddhanasasanam.
% Sabbapapassa akaranam
% Kusalass'upasampadil
% Sacittapariyodapanam
% Etam buddhanasasanam.
% Anupavado anupaghato
% Patimokkhe ca samvaro
% Mattannuta ca bhattasmim
% Pan tan ca sayanasanam
% Adhicitte ca ayogo
% Etam buddhanasasanan'ti.
% ```

% 原文对应的英文和泰语翻译。
% ```
%  * Khantī paramaṃ tapo titikkhā,
%    Patience is the supreme purifying practice,
%    (khuam-ot-thon khue khuam-ot-klan pen ta-ba yang-ying)
%  * nibbānaṃ paramaṃ vadanti buddhā,
%    Nibbāna is supreme, say the Buddhas.
%    (phra-phut-tha-jao tang-lai trat-wa phra-nip-phan pen-yiam)
%  * na hi pabbājito parūpaghātī samaṇo hoti paraṃ viheṭhayanto,
%    A mendicant does not harm others, A recluse oppresses no one.
%    (phu-lang-phlan phoo-uen mai-chue-wa pen ban-pha-chit, phoo- biat-bian phoo-uen mai-chue-wa pen-samana-loei)
%  * etaṃ buddhāna sāsanaṃ.
%    This is the teaching of all Buddhas.
%    (ni pen-kham-sorn khong phra-phut-tha-jao tang-lai)
%  * sabbapāpassa akaraṇaṃ,
%    Abstaining from all evil,
%    (kan-mai-tham-bap-puang nueng)
%  * kusalassūpasampadā,
%    Doing all good thing,
%    (kan-bam-phen ku-son hai-thueng-phrom nueng)
%  * sacittapariyodapanaṃ,
%    Cleansing of one's mind,
%    (kan-klan-chit khong-ton hai-chue-phaew nueng)
%  * etaṃ buddhāna sāsanaṃ.
%    This is the teaching of all Buddhas.
%    (ni pen-kham-sorn khong phra-phut-tha-jao tang-lai)
%  * anūpavādo,
%    Not insulting,
%    (kan-mai-khao-pai wa-rai-kan nueng)
%  * anūpaghāto,
%    Not harming,
%    (kan-kai-khao-pai lang-phlan-kan nueng)
%  * pātimokkhe ca saṃvaro,
%    Carefulness of the Fundamental Precepts,
%    (khuam-sam-ruam nai phra-pa-ti-mok nueng)
%  * mattaññutā ca bhattasmiṃ,
%    Moderating in food,
%    (khuam-pen phoo-ru-pra-man nai pho-cha-na-han nueng)
%  * pantañca sayanāsanaṃ,
%    Dwelling in solitude,
%    (kan-non kan-nang an-sa-ngad nueng)
%  * adhicitte ca āyogo,
%    Engaging in higher mental development,
%    (kan-pra-kop khuam-phian nai a-thi-jit nueng)
%  * etaṃ buddhāna sāsanaṃ.
%    This is the teaching of all Buddhas.
%    (ni pen-kham-sorn khong phra-phut-tha-jao tang-lai)
% ```

% 我需要的格式举例。和之前的不同，三种语言直接放在一起，方便用多语种吟诵。
% ```
% \section*{[12] ...}

% \begin{chantsubsection}
%     \chantnote
%     {\{Handa mayaṃ pattidāna gāthāyo bhaṇāma se\}}
%     \chantindent{\{Now let us recite the verses for sharing merit.\}}
% \end{chantsubsection}

% \begin{chantsubsection}
%     \chantnote
%     {Khanti paramam tapo titikkha}
%     Patience is the supreme purifying
%     \par
%     (Khuam-ot-thon khue khuam-ot-klan pen ta-ba yang-ying)

%     \chantnote
%     ...
% \end{chantsubsection}
% ```